\documentclass[12pt,a4paper]{article}
\usepackage[margin=2.5cm]{geometry}
\usepackage{hyperref}
\usepackage{parskip}
\usepackage{xcolor}
\usepackage{mdframed}
\usepackage{enumitem}
\usepackage{microtype}

% Reviewer comment box
\newmdenv[
  backgroundcolor=gray!10,
  linecolor=gray!40,
  linewidth=0.8pt,
  innerleftmargin=10pt,
  innerrightmargin=10pt,
  innertopmargin=8pt,
  innerbottommargin=8pt,
  skipabove=10pt,
  skipbelow=4pt
]{reviewbox}

% Response formatting
\newcommand{\reviewer}[1]{\medskip\noindent{\large\textbf{#1}}\par\smallskip}
\newcommand{\point}[1]{\medskip\noindent\textbf{Comment #1.}\enspace}
\newcommand{\response}{\noindent\textit{Response:}\enspace}
\newcommand{\change}{\noindent\textit{Change made:}\enspace}

\begin{document}

%----------------------------------------------------------------------
% Header
%----------------------------------------------------------------------
\begin{center}
  {\Large\textbf{Author's Response to Reviewers}}\\[6pt]
  {\large\textit{The Fibonacci Quarterly}}\\[4pt]
  {\normalsize Manuscript: ``A Time-Domain Analogue to Fibonacci Structure
  via Phase-Gated AR(2) Dynamics''}\\[4pt]
  {\normalsize Author: M.\ Whiteside \quad Date: \today}
\end{center}

\bigskip
\hrule
\bigskip

I thank the reviewer and the Academic Editor for their careful
reading of the manuscript and for the constructive and detailed feedback.
The revision addresses every point raised.
Below I respond to each comment in turn, quoting the reviewer's text
in shaded boxes, followed by my response and a precise description of
the change made.
All section numbers refer to the revised manuscript
(\textit{Whiteside\_FQ\_Revision\_v1}).

\bigskip
\hrule

%======================================================================
\reviewer{Reviewer 1}
%======================================================================

\point{1}

\begin{reviewbox}
It would strengthen the story for the reader if a clear rationale was
provided for analysis of data on rhythmic organoids as well as on tuft
cells, deep crypt secretory cells, and long-lived lineages.
\end{reviewbox}

\response
I agree that the original manuscript did not make the selection
rationale sufficiently explicit.
The PAR(2) framework requires three specific inputs: (i) a rhythmic
time-series from crypt-relevant tissue; (ii) a long-lived cell type
whose dynamics accumulate temporal perturbations across many renewal
cycles; and (iii) a niche-supporting cell type whose sustained
signalling activity can physically implement the PAR(2) memory kernels.
Each choice is now justified individually.
GSE157357 (Stokes et al.) is the only publicly available circadian
time-series that simultaneously profiles crypt stem cell behaviour in
wild-type and clock-disrupted conditions, making it uniquely suited to
test the PAR(2) kernel architecture.
Tuft cells are chosen because their lifespan ($\geq$28 days) far exceeds
the 3--5-day average crypt turnover, so they accumulate the effects of
eigenvalue perturbations across many cycles and provide the most
sensitive available readout.
DCS cells are chosen because their position at the crypt base---providing
continuous WNT, EGF, and Notch ligands to LGR5$^{+}$ stem cells---makes
them the natural physical implementors of the slow and fast memory
kernels $\alpha_2$ and $\alpha_1$ in the PAR(2) model.

\change
Section~5 now opens with a new structured paragraph giving these
three justifications explicitly before proceeding to the biological
discussion of each cell type.

%----------------------------------------------------------------------
\point{2}

\begin{reviewbox}
A potential weakness should be noted since the study by Boman et al.\
modeled human large intestinal crypts while the study by Stokes et al.\
involved data on crypts from mouse small intestinal organoids.
\end{reviewbox}

\response
This is a genuine and important limitation.
The species difference (human vs.\ mouse) and anatomical region
difference (large intestine vs.\ small intestine) mean that the AR(2)
fits cannot be taken as a direct quantitative test of Boman's
human-colon tissue code.
I have made this caveat much more prominent.

\change
A bold-labelled \textbf{Limitation: model--data mismatch} paragraph has
been added directly after the introduction of the Stokes et al.\ dataset
in Section~4.
The paragraph names the specific biological differences (LGR5$^{+}$
stem cell density, villus structure, crypt length, cancer
susceptibility) and states explicitly that these data serve as a
proof-of-concept demonstration only, not a validated quantitative test.
Direct validation against a human large-intestinal circadian time-series
is named as an important open objective.

\medskip
\noindent\textit{Note to reviewer.}
As reassurance that the Fibonacci-proximate AR(2) signal is not
species-specific, I note that preliminary analysis of the human
intestinal enteroid dataset (GSE161566, Rosselot et al.\ 2022 ---
the human subseries of GSE179028, 24 timepoints at 2-hour intervals)
shows consistent Fibonacci-proximate dynamics: a pre-specified
14-gene E-box target set clusters significantly closer to the $1/\varphi$
boundary in AR(2) parameter space than expression-matched random
gene sets ($p = 0.029$, permutation test).
This cross-species result is not included in the current revision,
which I have kept within the scope of the original submission,
but it supports the biological plausibility of the PAR(2) framework
in human intestinal tissue and will be developed in subsequent work.

%----------------------------------------------------------------------
\point{3}

\begin{reviewbox}
Reference~5 appears to contain a few inaccuracies regarding the author,
journal volume, and year. Specifically: the paper is by Edward
Bormashenko, not M.\ Pinto. It was published in 2022, but in a
different journal than the main \textit{Symmetry} journal (which was at
Volume~14 in 2022). It appears as Article~27 in Volume~2, Issue~3 of
the MDPI journal \textit{Foundations} (formerly part of the Symmetry
topic family).
\end{reviewbox}

\response
I thank the reviewer for the detailed correction.
The in-text attribution (``Bormashenko'') was already correct in the
submitted manuscript; the error resided in the reference list entry,
which incorrectly listed the author as M.\ Pinto.
The reviewer correctly identifies the article coordinates
(Vol.~2, Issue~3, Art.~27, 2022).
On the journal name: while the reviewer suggests \textit{Foundations},
the published DOI \texttt{10.3390/biophysica2030027} resolves
unambiguously to \textit{Biophysica}, a distinct MDPI journal.
The reference has been corrected accordingly.

\change
The reference list entry has been corrected: author changed from
``M.\ Pinto'' to ``Bormashenko, E.''; journal corrected from
\textit{Foundations} (as suggested by the reviewer) to
\textit{Biophysica} (per the published DOI);
volume, issue, and article number verified.
The corrected entry reads:
Bormashenko, E. Fibonacci sequences, symmetry and order in biological
patterns, their sources, information origin and the Landauer principle.
\textit{Biophysica.} 2(3):Art.\ 27, 2022.
doi:10.3390/biophysica2030027.

%----------------------------------------------------------------------
\point{4}

\begin{reviewbox}
To make the Fibonacci-consistent manifold more robust, perhaps a formal
definition or a visualization of this manifold could be provided. Where
exactly in the parameter plane does ``Boman-like'' behavior live?
\end{reviewbox}

\response
I agree that ``manifold'' was used without a precise definition.
The Fibonacci-consistent manifold is now defined concretely in
Section~6.

\change
A new paragraph at the end of Section~6 defines the manifold as the
subset of the stable AR(2) stationarity triangle---the biologically
admissible parameter space---in which the coefficient pair
$(\varphi_1,\varphi_2)$ lies close to the Fibonacci boundary landmark
at $(1,1)$, with the admissible neighbourhood bounded above and below
by the division-limit (Rule~3) and lifespan (Rule~4) constraints of
Boman's tissue code.
The previously vague future-work sentence (``formalise a
Fibonacci-consistent region in AR(2) parameter space'') has been removed
and replaced with a cross-reference to this definition, together with a
revised future-work sentence focused on experimental validation of the
manifold boundary rather than its creation.

%----------------------------------------------------------------------
\point{5}

\begin{reviewbox}
Can mapping the temporal parameters onto Boman's spatial rules be put
into context of the biology? In Table~1, the roots for genes are quite
small indicating stability. However, in a true Fibonacci recurrence the
roots involve the golden ratio ($+0.618\ldots$). While the data shows
stability (strongly damped oscillations), it is mathematically ``far''
from the unstable or marginally stable growth typically associated with
Fibonacci sequences. Perhaps it might be clarified how the heavily
damped regime relates to, or is consistent with, the Fibonacci analogy.
\end{reviewbox}

\response
This is an important conceptual point and I am grateful for it being
raised.
The Fibonacci recurrence $x_n = x_{n-1}+x_{n-2}$ has dominant root
$|\lambda|\approx 1.618$, which lies outside the AR(2) stationarity
triangle: exact Fibonacci dynamics would be explosive, not stable.
Biological tissues must operate with $|\lambda|<1$; the observed root
magnitudes in Table~1 range from 0.063 to 0.521, far from the explosive
boundary.
The biological content of the Fibonacci analogy is therefore that
homeostatic regulation holds AR(2) coefficients near the stable twin of
the Fibonacci eigenvalue, $1/\varphi\approx 0.618$, and Boman's five
rules are precisely the biological mechanisms that enforce this heavily
damped sub-unit-root regime.

\change
A ``Critical clarification'' paragraph has been added in Section~6
immediately after the five-rules mapping, making this distinction
explicit with reference to the measured eigenvalue magnitudes from
Table~1.

\bigskip
\hrule

%======================================================================
\reviewer{Reviewer 2 (Academic Editor)}
%======================================================================

\response
I thank the Academic Editor for the careful line-by-line review.
All formatting corrections have been applied.
I address each point below.

\point{1--12 (``e.g.'' without comma)}

\begin{reviewbox}
Throughout the manuscript, ``e.g.'' should be followed by a comma:
``e.g.,''.
\end{reviewbox}

\change
All instances of ``e.g.'' have been corrected to ``e.g.,''.
A \texttt{grep} search of the revised \texttt{.tex} source confirms that
zero instances of ``e.g.'' without a trailing comma remain.

%----------------------------------------------------------------------
\point{13 (Disclosure statement)}

\begin{reviewbox}
Add a Disclosure Statement reading ``No conflict of interest has been
reported by the author.''
\end{reviewbox}

\response
On reflection, a patent application relevant to the methodology
described in this work existed at the time of original submission and
must be declared to ensure full transparency.
UK patent application GB2518973.9, covering the AR(2) eigenvalue
methodology, was filed on 13 November 2025 --- three days before the
manuscript was submitted on 16 November 2025.
It should have been declared at submission; I correct this omission now.

\change
The Disclosure Statement has been added and updated to read:
``The author filed a UK patent application (GB2518973.9) related to the
AR(2) eigenvalue methodology described in this work on 13 November 2025,
three days prior to the original manuscript submission on 16 November 2025.
The author declares no other conflicts of interest.''

%----------------------------------------------------------------------
\point{14 (Reference formatting)}

\begin{reviewbox}
Rewrite the Reference section in the accepted Fibonacci Association
style: alphabetical order, Last Name First Initial format, no quotation
marks for article titles, abbreviated journal names.
\end{reviewbox}

\change
All 11 references have been reformatted to Fibonacci Quarterly style:
alphabetical order by first author's last name; Last Name, Initial(s)
format; no quotation marks around article titles; abbreviated journal
names (e.g., \textit{Fibonacci Quart.}, \textit{Biol.\ Cell},
\textit{Front.\ Immunol.}).

%----------------------------------------------------------------------
\point{15 (Phys.org news article)}

\begin{reviewbox}
Reference for the Phys.org news article must include an author name,
or the reference should be removed.
\end{reviewbox}

\change
The Phys.org news article reference has been removed from the reference
list entirely.

%----------------------------------------------------------------------
\point{16 (``and co-authors'' $\to$ ``et al.'')}

\begin{reviewbox}
Change ``and co-authors'' to ``et al.''
\end{reviewbox}

\change
The phrase ``Boman and co-authors'' has been changed to
``Boman et al.'' throughout.

%----------------------------------------------------------------------
\point{17 (Duplicate reference)}

\begin{reviewbox}
Remove the redundant reference, which is already included as another
reference.
\end{reviewbox}

\change
The duplicate reference to the Bormashenko \textit{Biophysica} paper
has been removed.
The reference list has been renumbered accordingly and now contains
11 unique entries.

%----------------------------------------------------------------------
\point{18 (Pages 11--19)}

\begin{reviewbox}
Are pages 11--19 extraneous and intended for a different document?
\end{reviewbox}

\response
Yes. Those pages belonged to supplementary material associated with the
original submission and were inadvertently included in the submitted PDF.
They are not part of the reply and have been removed.

\change
The revised manuscript is 8 pages in length.

\bigskip
\hrule
\bigskip

\noindent\textbf{Summary of all changes.}
The table below provides a concise cross-reference.

\medskip
\begin{center}
\small
\begin{tabular}{lll}
\hline
\textbf{Reviewer} & \textbf{Comment} & \textbf{Location in revision} \\
\hline
R1 & 1. Cell-type rationale & Section~5, opening paragraph \\
R1 & 2. Mouse/human mismatch & Section~4, bold Limitation paragraph \\
R1 & 3. Reference attribution & Reference~[3]; in-text citation \\
R1 & 4. Manifold definition & Section~6, final paragraph \\
R1 & 5. Damped regime clarification & Section~6, Critical clarification paragraph \\
\hline
Editor & 1--12. e.g.\ $\to$ e.g., & Throughout \\
Editor & 13. Disclosure statement & AI Use Declaration / Disclosure section \\
Editor & 14. Reference formatting & Full reference list \\
Editor & 15. Phys.org reference & Removed \\
Editor & 16. et al. & Section~1 \\
Editor & 17. Duplicate reference & Reference list (11 unique entries) \\
Editor & 18. Pages 11--19 & Removed; revised MS is 8 pages \\
\hline
\end{tabular}
\end{center}

\bigskip
I hope that the revised manuscript now fully meets the standards of
\textit{The Fibonacci Quarterly} and am grateful for the opportunity
to improve the work.

\end{document}
